\section{Diffuse Heisenberg symbols and voltage approximation}
\label{sec:analysis}

We assign a Heisenberg voltage and a bounded matrix weight to every surviving
old edge.  The weighted incidences will approximate prescribed compact Bott
symbols in the \emph{full} regular operator norm, while the voltages will
separate all the short closed walks specified below.  These requirements use
different features of the same distribution: its weighted mean realizes the
symbol, and its diffuse horizontal coordinates prevent trivial holonomy.
The horizontal coordinates remain diffuse when the matrix weight is zero.

For the voltage estimates we fix a surviving tower furnished by
Proposition~\ref{prop:pruning}, with any prescribed finite collection of
incidence types.  All subsequent edge draws are independent on this fixed
graph; the sequential table law has already served to select its incidences.
Before deletion, a type between an
\(\mathscr A_i\)-layer and a \(\mathscr B_k\)-layer, \(\abs{i-k}\leq1\), is
biregular.  Denote its degrees by \(d_{\tau,ik}\) and \(D_{\tau,ik}\), its
normalizing factor by
\[
 q_{\tau,ik}=(d_{\tau,ik}D_{\tau,ik})^{1/2},
\]
and its surviving normalized incidence by \(U_{\tau,ik}\).  All normalizations
use the original degrees.  For fixed types and fixed \(b\), there are positive
constants \(c_*,C_*\), independent of height and of sufficiently large \(s\),
such that
\begin{equation}
\label{eq:analysis-degree-bounds}
 c_*d\leq d_{\tau,ik}\leq C_*d,\qquad
 c_*D\leq D_{\tau,ik}\leq C_*D,\qquad
 c_*q\leq q_{\tau,ik}\leq C_*q.
\end{equation}
In particular \(\norm{U_{\tau,ik}}\leq1\), because it is a compression of a
normalized biregular incidence.  The full surviving base graph has maximum
degree at most \(C_*D\).  We shall use its base-girth bound \(>.93m_i\) on
comparable layers, and its local absence of Heisenberg laws of length at most
\(6.04m_i\) based at height i.  These are inputs from
Proposition~\ref{prop:pruning}; its proof applies
Lemma~\ref{lem:heisenberg-law} with the appropriate local scale comparison.
The base-girth bound will control the centered moments, and the nonlaw bound
will control the probability of trivial holonomy.

For definiteness the formulas are written in the left regular representation
\(\lambda\) of \(\HH\).  The inversion unitary
\(\mathcal I\delta_g=\delta_{g^{-1}}\) carries \(\lambda(a)\) to the right
regular representation \(\rho(a)\delta_g=\delta_{ga^{-1}}\).  Thus all statements,
including the matrix-valued ones, hold in either convention.  The voltage
and incidence orientations are fixed as follows.  A positive old edge with
label \(\sigma\) runs from \(A_v\) to \(B_w\), and its forward voltage map is
\((v,f)\mapsto(w,fh_e)\).  The block \(U_{\tau,ik}\) maps B-coordinates to
A-coordinates, so its deck entry acts by
\[
 \rho(h_e)\delta_f=\delta_{fh_e^{-1}}.
\]
Thus the sampled convolution variable \(v_e\) below is assigned as
\(h_e=v_e\); the reverse block is its adjoint, with inverse voltage.
Equation~\eqref{eq:assembly-edge-direction} later identifies this same
deck action in the embedded Cayley graph.

\subsection{Compact Bott representatives and fixed smooth families}

We use a smooth compactly supported form of the classical Bott representative
associated to the Hopf line bundle; compare \cite[Section~2]{AB64}.
The following formula fixes its support and orientation conventions.  Put
\[
 h(t)=\begin{cases}e^{-1/t},&t>0,\\0,&t\leq0,\end{cases}
 \qquad
 \zeta(r)=\frac{h(r-1/4)}{h(r-1/4)+h(1-r)},\qquad r\geq0,
\]
and set \(\theta(r)=2r+(\pi-2r)\zeta(r)\).  Thus \(\theta(r)=2r\) near zero and
\(\theta(r)=\pi\) for \(r\geq1\), with a smooth constant extension at the latter
endpoint.  For \(r=(t^2+u^2)^{1/2}>0\), define
\begin{equation}
\label{eq:compact-bott-profile}
 p(t,u)=\frac12
 \begin{pmatrix}
 1+\cos\theta(r)&\dfrac{\sin\theta(r)}{r}(t-iu)\\[3pt]
 \dfrac{\sin\theta(r)}{r}(t+iu)&1-\cos\theta(r)
 \end{pmatrix},
 \qquad e_\infty=\begin{pmatrix}0&0\\0&1\end{pmatrix}.
\end{equation}
At zero use the continuous extension.  The functions
\(\sin(2r)/r\) and \(\cos(2r)\) are analytic functions of \(r^2\) near zero;
hence this extension is smooth.  At \(r=1\) the difference from the constant
matrix is flat.  Direct multiplication, or the identity
\((n\cdot\sigma)^2=1\) for the unit vector
\[
 n=(\sin\theta\cos\phi,\sin\theta\sin\phi,\cos\theta),
\]
shows that \(p=p^*=p^2\) and \(\operatorname{Tr}p=1\).  In particular
\(p-e_\infty\) is a smooth compactly supported matrix function.

The compactification of this projection has degree of absolute value one.
Indeed, its local unit eigenvectors are
\[
 \binom{\cos(\theta/2)}{e^{i\phi}\sin(\theta/2)}
 \quad\hbox{near the origin},\qquad
 \binom{e^{-i\phi}\cos(\theta/2)}{\sin(\theta/2)}
 \quad\hbox{near the exterior constant region}.
\]
On their overlap the second is \(e^{-i\phi}\) times the first, so the clutching
function has winding number \(-1\).  Equivalently, with the corresponding
choice of Chern-form sign,
\[
 \frac{1}{2\pi i}\int_{\mathbb R^2}
       \operatorname{Tr}(p\,dp\wedge dp)
 =\frac12\int_0^1\theta'(r)\sin\theta(r)\,dr=1.
\]
The winding calculation is independent of this sign convention.  Thus the
rank-zero difference represents one of the two Bott generators.

Fix
\[
 \alpha=\frac{\sqrt5-1}{2},\qquad
 N_i=2^{200H_i},\quad Z_i=2^{500H_i},\quad J_i=2^{50H_i}.
\]
The commuting pair \((z,x)\) identifies its group algebra with the continuous
functions of torus coordinates \((\vartheta,\xi)\).  In fixed coordinate
charts around \((\alpha,0)\), set
\begin{equation}
\label{eq:scaled-bott-projection}
 \mathsf p_i(\vartheta,\xi)
   =p\bigl(Z_i(\vartheta-\alpha),N_i\xi\bigr),
 \qquad \beta_i=\mathsf p_i-e_\infty,
\end{equation}
and extend \(\mathsf p_i\) by \(e_\infty\).  The initial height is chosen so
that every support lies inside the charts.  A change of variables in the
preceding integral, or the same clutching function, proves that every
\(\beta_i\) is the rank-zero Bott generator up to the fixed orientation sign.

Adjacent representatives are joined by the explicit rescaling
\[
 p\bigl(Z_i^{1-v}Z_k^v(\vartheta-\alpha),
        N_i^{1-v}N_k^v\xi\bigr),\qquad 0\leq v\leq1,
 \quad\abs{i-k}=1.
\]
Since \(H_k-H_i=\pm b\), the ratios of these scales are fixed.  In rescaled
coordinates this is a compact smooth family, with all derivatives bounded
independently of height.  The matrix rotations
\[
 R(v)=\begin{pmatrix}\cos(\pi v/2)&-\sin(\pi v/2)\\
                      \sin(\pi v/2)& \cos(\pi v/2)\end{pmatrix}
       \otimes I_2
\]
act by conjugation on \(\mathsf p_i\oplus e_\infty\), joining it to
\(e_\infty\oplus\mathsf p_i\).  For sampling we subtract the constant
projection \(e_\infty\oplus e_\infty\) from this rotation family and
\(e_\infty\) from the rescaling family.  Entries of these compact
differences, their rescalings, and smooth cutoffs equal to one on
their supports belong to a fixed compact smooth profile family.  One may
choose all cutoffs inside rectangles larger by a constant depending only on
\(b\).  In particular, the carrier at height \(i+1\) may contain the supports
of the height-\(i\) profiles used in the adjacent homotopy.  Throughout the
final accuracy sequence we keep
\(b\) fixed, so these profile bounds and moduli of continuity are uniform also
in the construction index.

More generally we will use matrix-valued functions of the form
\begin{equation}
\label{eq:profile-family}
 f_i(\vartheta,\xi)=F\bigl(Z_i(\vartheta-\alpha),N_i\xi\bigr),
\end{equation}
where \(F\) ranges over a fixed compact family in
\(C_c^\infty(\mathbb R^2;M_\ell(\mathbb C))\), supported in one fixed
rectangle.  Zero functions and any prescribed height selections are allowed.
The matrix dimension \(\ell\) is fixed.  Passing between adjacent reference
heights merely enlarges this profile family by fixed rescalings.

\subsection{Spreading on disjoint torus translates}

We record explicitly the separation property used here.  If \(p\) is the nearest
integer to \(n\alpha\), and
\(\alpha'=(-\sqrt5-1)/2\), then
\[
 (n\alpha-p)(n\alpha'-p)=p^2+np-n^2
\]
is a nonzero integer.  Moreover
\(\abs{n\alpha'-p}\leq\sqrt5\abs n+1/2\leq3\abs n\) for \(n\ne0\).
Consequently
\begin{equation}
\label{eq:golden-separation}
 \norm{n\alpha}_{\mathbb R/\mathbb Z}\geq\frac{1}{3\abs n},
 \qquad n\ne0.
\end{equation}

Choose a nonzero real function \(a\in C_c^\infty((-1,1))\), and put
\[
 a_{ij}=\frac{a(j/J_i)}{\bigl(\sum_{k\in\mathbb Z}
                      \abs{a(k/J_i)}^2\bigr)^{1/2}},\qquad
 T_i=\sum_{j\in\mathbb Z}a_{ij}\lambda(y)^j.
\]
All sums defining \(T_i\) are finite and \(\sum_j\abs{a_{ij}}^2=1\).
Riemann sums show that the denominator is bounded above and below by positive
constant multiples of \(J_i^{1/2}\), once the initial height is large.

Let \(E_i\) be the Borel spectral projection of an enlarged rectangle
\[
 \abs{\vartheta-\alpha}\leq R/Z_i,\qquad
 \abs\xi\leq R/N_i,
\]
large enough to contain all symbols and cutoffs incident to height \(i\).
These Borel projections are auxiliary support operators; they are not asserted
to belong to the continuous group algebra.

\begin{lemma}[Supported spreading]
\label{lem:supported-spreading}
For sufficiently large initial height, the family of projections
\[
 \lambda(y)^jE_i\lambda(y)^{-j},\qquad \abs j\leq J_i,
\]
is pairwise orthogonal.  Consequently
\[
 E_iT_i^*T_iE_i=E_i,\qquad P_i=T_iE_iT_i^*
\]
is a projection, and
\begin{equation}
\label{eq:spread-trace}
 \tau_{\HH}(P_i)=\tau_{\HH}(E_i)
 \leq C_R2^{-700H_i}.
\end{equation}
If \(f\) is supported in both enlarged rectangles at heights \(i\) and \(k\),
and \(g\) in both enlarged rectangles at heights \(k\) and \(j\), then
\[
 (T_i fT_k^*)(T_k gT_j^*)=T_i(fg)T_j^*,\qquad
 (T_i fT_k^*)^*=T_k f^*T_i^*,\qquad
 \norm{T_i fT_k^*}\leq\norm f.
\]
\end{lemma}

\begin{proof}
The relation \([x,y]=z\) gives \(y^jxy^{-j}=xz^{-j}\), so conjugating the
torus rectangle changes the \(\xi\)-coordinate by \(-j\vartheta\).  If
\(0<\abs{j-k}\leq2J_i\) and \(\abs{\vartheta-\alpha}\leq R/Z_i\),
\[
 \norm{(j-k)\vartheta}_{\mathbb R/\mathbb Z}
 \geq\frac{1}{6J_i}-2RJ_i/Z_i.
\]
This exceeds \(2R/N_i\) at all sufficiently large heights because
\(N_i=J_i^4\) and \(Z_i=J_i^{10}\).  The translated supports are disjoint.
Expansion of \(E_iT_i^*T_iE_i\) now leaves only its diagonal terms, proving
the first identity.  All subsequent product and norm identities follow by
inserting the auxiliary support projections.  The trace on the torus algebra
is Haar measure: it vanishes on every nonidentity character
\(\lambda(z)^c\lambda(x)^n\), which determines its integral on continuous
functions and hence on Borel spectral projections.  The trace property and
\(E_iT_i^*T_iE_i=E_i\) give \(\tau(P_i)=\tau(E_i)\), and the rectangle has
area at most \(4R^2/(Z_iN_i)\).  In particular the number of translates does
not multiply the trace of this isometric image.
\end{proof}

\subsection{The simultaneous approximation statement}

The spread symbols now have the support and product identities needed for the
Bott construction.  We next realize any finite prescribed family of them by
weighted graph incidences, while choosing the same voltages to eliminate
short closed lifts.

\begin{proposition}[Simultaneously sampled symbols]
\label{prop:sampled-symbols}
Fix b, a finite matrix dimension, a positive accuracy \(\eta\), and a finite
list of compact smooth symbol profiles and height selections as above.  The
list may include cutoffs, both parities, the initial-height selection, and any
finite mesh of the adjacent rescaling and rotation families.  Allocate a
prescribed incidence type for every required same or adjacent layer symbol.
For sufficiently large s and then sufficiently large initial height, the
surviving base graph admits a voltage on every edge and finite-net matrix weights
such that the following hold simultaneously.
\begin{enumerate}
\item The voltage cover has no nontrivial cyclically reduced closed walk of
length at most \(6.03m_i\) based at an old vertex of height i.
\item All weighted blocks may use the common denominator q.  If
\(\mathcal V_{\tau,ik}\) is that normalized weighted convolution block and
the prescribed symbol is \(S_{\tau,ik}=T_if_{\tau,ik}T_k^*\), then
\[
 \norm{\mathcal V_{\tau,ik}
           -U_{\tau,ik}\otimes S_{\tau,ik}}\leq\eta
\]
in the full regular operator norm, at every height.  The tensor notation here
means that each scalar incidence entry multiplies the same matrix-valued deck
convolution symbol.
\item Weights are read from one finite set, independent of height.  Appending
their net indices to the existing edge labels preserves the labelled-graph
uniqueness properties and gives
\[
 \log_2\abs{\mathcal A}
 \leq1.00005s+O(\log s)+O_{b,\eta,\mathrm{profiles}}(1)
 <\log_2(q^2)
\]
for sufficiently large s.  In particular \(\sqrt{\abs{\mathcal A}}/q\)
tends to zero as s increases after the fixed data.  Zero-weight edges are
retained in the voltage cover.
\end{enumerate}
\end{proposition}

The proof is completed in Section~\ref{subsec:analysis-simultaneous}.
It has three parts.  A Fourier coefficient envelope will produce bounded
matrix weights with the desired mean and independent diffuse horizontal
coordinates.  A polynomial zero estimate will then control the holonomy of
each short nonlaw.  Finally, a finite regular-norm test and a centered moment
count will control the fluctuations of every weighted incidence block.
Their failure probabilities are summable over the height, so one choice
satisfies all the requirements.

\subsection{A product envelope for the Fourier coefficients}

To sample the two horizontal coordinates independently with small atoms, we
need a coefficient bound that separates their contributions.  The envelope
below is uniform over the \(O(J_i)\) possible y-powers; the sum of its
x-frequency majorant is \(O(J_i^{-1})\).  These two factors compensate,
giving both a uniform total coefficient mass and the product law used in
Lemma~\ref{lem:product-sampler}.  Control of the central Fourier tails will
make every edge distribution finite.

Every element of \(\HH\) has a unique form \(x^ny^lz^c\).  For a matrix
Fourier series, write
\[
 A=\sum_{n,l,c}A_{n,l,c}\lambda(x^ny^lz^c),\qquad
 \norm A_1=\sum_{n,l,c}\norm{A_{n,l,c}}.
\]
This one-norm dominates the
regular operator norm by the triangle inequality.

We first state an elementary rescaling estimate.  If \(g\) is a smooth compactly
supported matrix function on \(\mathbb R\), periodized as
\(g(Z(\vartheta-\alpha))\) inside one chart, its central Fourier coefficients
are
\[
 c_k=Z^{-1}e^{-2\pi i k\alpha}\widehat g(k/Z).
\]
For every nonnegative integer \(A\), integration by parts \(A+2\) times,
together with the trivial integral bound for \(\abs{k/Z}\leq1\), gives
\begin{equation}
\label{eq:rescaled-fourier-estimate}
 \sum_k(1+\abs k/Z)^A\norm{c_k}
 \leq C_A\sum_{j=0}^{A+2}\norm{g^{(j)}}_{L^1(\mathbb R)}.
\end{equation}
Indeed the pointwise bound after those integrations is a constant times
\((1+\abs{k/Z})^{-A-2}\) times the displayed derivative sum, and
\(Z^{-1}\sum_k(1+\abs{k/Z})^{-2}\) is bounded uniformly for \(Z\geq1\).
The same argument uses operator norms and therefore applies to a fixed matrix
amplification without any scalar-entry assumption.

\begin{lemma}[Fourier envelope]
\label{lem:fourier-envelope}
Let \(\abs{i-k}\leq1\), and let \(f\) belong to the profile family
\eqref{eq:profile-family}, expressed at reference height \(i\).  Write
\[
 T_i fT_k^*=\sum_{n,l,c}A_{n,l,c}\lambda(x^ny^lz^c),
 \qquad N=N_i,\quad Z=Z_i,\quad J=J_i.
\]
For every pair of nonnegative integers \(A,P\), with \(P\geq3\),
\begin{equation}
\label{eq:weighted-central-envelope}
 \sum_c(1+\abs c/Z)^A\norm{A_{n,l,c}}
 \leq C_{A,P}N^{-1}(1+\abs n/N)^{-P}
       (1+J\norm{n\alpha}_{\mathbb R/\mathbb Z})^{-3},
 \qquad \abs l\leq C_bJ,
\end{equation}
and the coefficients vanish for larger \(\abs l\).  The constants are uniform
in height and in the compact profile family.  In particular these symbols
have a uniformly bounded Fourier one-norm.  For each \(\eta>0\) they can be
truncated, with one-norm error at most \(\eta\), to
\begin{equation}
\label{eq:fourier-box}
 \abs n\leq L_\eta N,\qquad \abs l\leq C_bJ,\qquad
 \abs c\leq L_\eta Z,
\end{equation}
where \(L_\eta\) is independent of height.
\end{lemma}

\begin{proof}
We first bound the central coefficients at fixed n and l, and then sum the
resulting horizontal majorant.
Fourier transform only in \(\xi\).  With \(t=Z(\vartheta-\alpha)\), its
\(n\)-th coefficient is
\[
 f_n(\alpha+t/Z)=N^{-1}\Phi(t,n/N),\qquad
 \Phi(t,v)=\int_{\mathbb R}F(t,u)e^{-2\pi ivu}\,du.
\]
The function \(\Phi\) is supported in one fixed compact t-interval and is
Schwartz in v, uniformly with all t-derivatives and all profile parameters.
For example integration by parts in u bounds
\(\norm{\partial_t^a\Phi(t,v)}\) by \(C_{a,M}(1+\abs v)^{-M}\) for every
fixed \(a,M\).

The Heisenberg commutation formula gives the coefficient at y-power l, before
central Fourier transformation, as
\[
 f_n(\vartheta)S_{n,l}(\vartheta),\qquad
 S_{n,l}(\vartheta)=\sum_j b_{j,l}e^{-2\pi i n j\vartheta},\qquad
 b_{j,l}=a_{ij}\overline{a_{k,j-l}}.
\]
The ratios \(J_k/J_i\) are fixed.  Thus \(b_{j,l}\) is supported on
\(O_b(J)\) indices, is zero unless \(\abs l\leq C_bJ\), and its r-th discrete
differences satisfy
\[
 \sum_j\abs{\Delta^r b_{j,l}}\leq C_rJ^{-r}.
\]
These statements follow by applying the fundamental theorem of calculus r
times to the two fixed bump factors, whose normalization is comparable to J.

It is essential to retain the central residual phase.  Treat
\(e^{-2\pi i n jt/Z}\) as part of the amplitude and sum by parts against
\(e^{-2\pi i n j\alpha}\).  For all fixed nonnegative integers r,a, the
discrete product rule and \(\abs j\leq C_bJ\) give
\[
 \sum_j\left|\Delta^r\partial_t^a
              \bigl(b_{j,l}e^{-2\pi i n jt/Z}\bigr)\right|
 \leq C_{r,a}J^{-r}(1+\abs n J/Z)^{r+a}.
\]
The trivial sum bound and r summations by parts therefore imply
\begin{equation}
\label{eq:central-residual-derivatives}
 \left|\partial_t^aS_{n,l}(\alpha+t/Z)\right|
 \leq C_{r,a}(1+\abs n J/Z)^{r+a}
          (1+J\norm{n\alpha}_{\mathbb R/\mathbb Z})^{-r}.
\end{equation}
This also covers \(n=0\) by the trivial bound.  To see the factor
\(1+J\norm{n\alpha}\) precisely, summation by parts divides by
\(\abs{1-e^{2\pi in\alpha}}^r\), which is bounded below by a constant times
\(\norm{n\alpha}^r\); taking the minimum with the trivial bound gives the
stated expression.

Apply \eqref{eq:rescaled-fourier-estimate} to
\(g(t)=N^{-1}\Phi(t,n/N)S_{n,l}(\alpha+t/Z)\), using r=3 in
\eqref{eq:central-residual-derivatives}.  Since
\[
 NJ/Z=2^{-250H_i}\leq1,\qquad
 1+\abs n J/Z\leq1+\abs n/N,
\]
the polynomial factors introduced by every fixed number of central derivatives
are absorbed by choosing the Schwartz exponent of \(\Phi\) larger.  This proves
\eqref{eq:weighted-central-envelope}, for arbitrarily large n as well as for
the principal x-frequency range.

It remains to obtain the compensating factor \(J^{-1}\) when summing over n.
On any interval of M
consecutive integers the points \(n\alpha\bmod1\) are separated by at least
\(1/(3M)\), by \eqref{eq:golden-separation}.  Hence at most \(1+6Mu\) lie at
distance at most u from any fixed point of the circle.  Integrating the
identity
\[
 (1+Jv)^{-3}=\int_v^\infty3J(1+Ju)^{-4}\,du
\]
against this counting bound gives
\begin{equation}
\label{eq:rotation-packing}
 \sum_{n\in I}(1+J\norm{n\alpha})^{-3}
 \leq C(1+M/J).
\end{equation}
For u beyond half the circle the same upper bound still bounds the count, so
the integral is valid as written.  Split the integers into the central interval
\(\abs n\leq N\) and dyadic annuli of lengths comparable to \(2^vN\).
Since \(N\geq J\), \eqref{eq:rotation-packing} yields
\begin{equation}
\label{eq:horizontal-envelope-sum}
 \sum_n N^{-1}(1+\abs n/N)^{-3}
              (1+J\norm{n\alpha})^{-3}\leq C/J.
\end{equation}
Summation over the \(O_b(J)\) possible l proves the one-norm assertion.  Using
a larger exponent P gives arbitrarily rapid polynomial tails when
\(\abs n>L N\); using a larger A gives the same statement for
\(\abs c>L Z\).  Their sum proves \eqref{eq:fourier-box}.
\end{proof}

\subsection{Diffuse product sampling with bounded matrix weights}

We turn the envelope into a probability distribution.  Its independent
horizontal marginals carry enough mass to dominate every Fourier coefficient;
the remaining coefficient is placed in a uniformly bounded matrix weight.
Unused mass is retained in the horizontal law even where the weight is zero.

\begin{lemma}[Finite product sampler]
\label{lem:product-sampler}
For every fixed \(\eta>0\), each symbol in Lemma~\ref{lem:fourier-envelope}
has a finite distribution of pairs \((W,v)\), with
\(W\in M_\ell(\mathbb C)\) and \(v=x^ny^lz^c\in\HH\), satisfying
\[
 \norm W\leq M_\eta,\qquad
 \norm{\mathbb E[W\lambda(v)]-T_ifT_k^*}_1\leq\eta.
\]
The horizontal coordinates n and l are independent.  Their maximal atoms
satisfy
\begin{equation}
\label{eq:horizontal-atoms}
 \sup_n\Pr(n)=O_b(J_i/N_i),\qquad
 \sup_l\Pr(l)=O_b(J_i^{-1}).
\end{equation}
All possible group coordinates, including those of inverses, are bounded by
\(2^{a_0H_i}\) for a fixed exponent \(a_0\), for example \(a_0=502\), after
increasing the initial height.  The weights can additionally be restricted to
one finite adjoint-closed net, at the expense of any prescribed further
one-norm error.  All constants and the net are independent of height.
\end{lemma}

\begin{proof}
Truncate the Fourier series to \eqref{eq:fourier-box}, with error at most
\(\eta/2\), and denote its matrix coefficients by \(A_{n,l,c}\).  Set
\[
 w_n=N_i^{-1}(1+\abs n/N_i)^{-3}
             (1+J_i\norm{n\alpha})^{-3}.
\]
By Lemma~\ref{lem:fourier-envelope},
\(R_{n,l}:=\sum_c\norm{A_{n,l,c}}\leq C_1w_n\), while
\(\sum_nw_n\leq C_0/J_i\).  On the finite allowed n-interval begin a
probability measure with masses \(J_iw_n/(2C_0)\), whose total is at most
one half.  Distribute the remaining mass uniformly on
\(\{-N_i,\ldots,N_i\}\).  Call the result \(\mu_i\).  Let \(\nu_i\) be
uniform on the integer interval containing all allowed l, of length comparable
to \(J_i\).  Then \eqref{eq:horizontal-atoms} holds and
\[
 R_{n,l}\leq M_\eta\mu_i(n)\nu_i(l)
\]
for a height-independent constant.  The finite cutoff L and fixed neighboring
ratios are harmless in this inequality.

Draw n and l independently from \(\mu_i\) and \(\nu_i\).  If \(R_{n,l}>0\),
draw c conditionally with probabilities \(\norm{A_{n,l,c}}/R_{n,l}\) on the
nonzero coefficients, and set
\[
 W=\frac{R_{n,l}}{\mu_i(n)\nu_i(l)}
         \frac{A_{n,l,c}}{\norm{A_{n,l,c}}}.
\]
If \(R_{n,l}=0\), set \(W=0\) and choose \(c=0\).  Thus zero-weight draws
still have exactly the same diffuse horizontal product law.  Directly summing
the expectation recovers every truncated coefficient, and \(\norm W\leq
M_\eta\).

Choose a finite net in the closed matrix ball of radius \(M_\eta\), include
zero, close it under adjoint, and round W within \(\eta/2\), keeping zero
fixed.  The resulting expectation changes in one-norm by at most
\(\mathbb E\norm{W-W_{\rm rounded}}\leq\eta/2\).  For a reversed edge use
the adjoint weight and inverse voltage rather than an independent draw.
The net cardinality depends on the accuracy and the fixed matrix dimension,
not on height.

The coordinate bounds follow from \eqref{eq:fourier-box}.  In the chosen normal
form,
\[
 (n,l,c)(n',l',c')=(n+n',l+l',c+c'-ln'),\qquad
 (n,l,c)^{-1}=(-n,-l,-c-ln).
\]
The product \(\abs{nl}\) is bounded by a constant times \(2^{250H_i}\), much
smaller than \(2^{500H_i}\).  Every fixed prefactor is absorbed by raising
the initial height, giving the stated exponent for both voltages and inverses.
\end{proof}

The zero symbol is included in this construction.  In particular a type may
have mean zero on any prescribed set of heights, such as all but the initial
height or one parity.  This requires no new height labels: the finite weight
net already contains zero.  The underlying voltage edge is retained.

\begin{lemma}[Polynomial avoidance]
\label{lem:polynomial-avoidance}
Use independent edge draws as in Lemma~\ref{lem:product-sampler}.
For a fixed base walk, reuse its one voltage
variable on every occurrence of an unoriented edge and use its inverse in the
reverse direction.  If the associated word is not a law in \(\HH\), the
probability that its holonomy is the identity is at most
\[
 C_b2^{-50H_{\min}},
\]
where \(H_{\min}\) is the smallest reference height of an edge in the walk.
The estimate is unaffected by matrix weights, including zero weights.
\end{lemma}

\begin{proof}
The polynomial estimate is the maximal-atom version of the usual polynomial
zero bound \cite[Lemma~1]{Schwartz80}; we include its short induction.
If independent real-valued variables have atoms at most \(\kappa\), a
nonzero polynomial of total degree d vanishes with probability at most
\(d\kappa\).  Induct on the number of variables.  If its degree in the last
variable is s, its leading coefficient is a nonzero polynomial in the preceding
variables of degree at most \(d-s\).  By induction that coefficient vanishes
with probability at most \((d-s)\kappa\); otherwise the one-variable polynomial
has at most s roots and costs at most \(s\kappa\).  Variables absent from the
polynomial are discarded.  A nonzero constant has zero vanishing probability,
starting the induction.

By the normal-form multiplication law just displayed, the horizontal
coordinates of any word are integral linear combinations of the edge
coordinates \(n_e,l_e\), with coefficients equal to signed occurrence counts.
If any such count is nonzero, one horizontal polynomial is nonzero.  If all
counts vanish, every central coordinate \(c_e\) cancels, and the central
coordinate of the word is a polynomial of degree at most two in the
\(n_e,l_e\).  Were that polynomial identically zero, the word would vanish for
every assignment in \(\HH\), contrary to the hypothesis.  Thus a nonzero
polynomial of degree at most two is tested on independent horizontal variables.
Their maximal atoms are bounded by \(C_b2^{-50H_{\min}}\), since
\(J_i/N_i=2^{-150H_i}\).  Apply the polynomial bound.  No conditioning on the
central coordinate or on the matrix weight is performed.  This proves the
assertion also for an edge whose realized weight is zero.
\end{proof}

\subsection{Finite tests for the full regular norm}

A sampled block acts on a finite base-layer space tensored with
\(\ell^2(\HH)\).  To control its full operator norm, we first test it in a
finite Heisenberg box, with exponentially small norm error and dimension at
most exponential in the height scale.  The subsequent moment order,
proportional to \(H_i/\log_2 d\), will overcome this dimension factor.

\begin{lemma}[Heisenberg boxes and regular norm]
\label{lem:heisenberg-folner-test}
Suppose a finite set of group elements and their inverses have coordinates
bounded by \(2^{a_0H}\), with \(H\geq1\).  For every fixed B there is a finite
box \(\mathcal B_H\subset\HH\), of cardinality at most \(2^{K_0H}\), such that
\[
 \frac{\abs{\mathcal B_H\triangle v\mathcal B_H}}
      {\abs{\mathcal B_H}}\leq C2^{-BH}
\]
for all these elements.  The exponent \(K_0\) depends on \(a_0,B\), but not
on the number of elements.

Let $A$ be a selfadjoint finite matrix of regular convolution sums whose
convolution support is contained in the finite set just specified, with any
fixed matrix amplification.  Let R bound its absolute coefficient row sums;
by selfadjointness it also bounds the column sums.  If \(Q_H\) retains the
deck coordinates in \(\mathcal B_H\), then
\begin{equation}
\label{eq:full-regular-norm-test}
 \norm{Q_HAQ_H}\leq\norm A
 \leq\norm{Q_HAQ_H}+C2^{-BH}R.
\end{equation}
The same statement holds for rectangular matrices after selfadjoint
off-diagonal doubling, using the larger row or column coefficient bound.
\end{lemma}

\begin{proof}
For an integer T let
\[
 \mathcal B(T)=\{x^ny^lz^c:\abs n,\abs l\leq T, \abs c\leq T^2\}.
\]
If \(v\) has coordinates bounded by \(R_0\) and \(T\geq R_0\), left or right
multiplication by \(v\) changes a horizontal coordinate by at most \(R_0\) and the
central coordinate by at most \(R_0+R_0T\).  This follows directly from the
normal-form multiplication law.  Points outside the corresponding boundary
strips remain in the box.  The total strip fraction is at most \(C R_0/T\).
Multiplication is a bijection, so the same bound, with a factor two, controls
the symmetric difference.  Take
\(T=\lceil2^{(a_0+B+2)H}\rceil\).  The box has
\((2T+1)^2(2T^2+1)=2^{O(H)}\) elements, proving the first assertion uniformly
over all possible translations.

For the norm estimate, let \(Q_t\) retain the deck coordinates in the right
translate \(\mathcal B_Ht\).  Right translation commutes with left convolution,
so every \(Q_tAQ_t\), acting on its range, has the same norm as \(Q_HAQ_H\).
For a finitely supported vector \(\xi\),
\[
 \sum_t\norm{Q_t\xi}^2=\abs{\mathcal B_H}\norm\xi^2.
\]
Average \(\ip{AQ_t\xi}{Q_t\xi}\) over t and divide by
\(\abs{\mathcal B_H}\).  In this average the coefficient of \(\lambda(v)\)
is multiplied by \(\abs{\mathcal B_H\cap v\mathcal B_H}/\abs{\mathcal B_H}\),
which differs from one by at most \(C2^{-BH}\).  The averaged quadratic form
has absolute value at most \(\norm{Q_HAQ_H}\norm\xi^2\).  The difference
from A has norm at most \(C2^{-BH}R\): to verify this last assertion, apply
Cauchy--Schwarz to a matrix of nonnegative entry-norm bounds whose row and
column sums are at most that number.  Density and the selfadjoint quadratic
form formula prove \eqref{eq:full-regular-norm-test}.  None of these estimates
separates coordinates of the matrix amplification.  Doubling a rectangular
matrix gives the final statement.
\end{proof}

\subsection{The centered moment count uses base girth}

We use the classical centered-trace method: independence eliminates walks
containing a singleton edge, and the remaining closed walks bound a high
moment of the norm; compare \cite[Sections~3.1--3.3]{FK81}.  Here base girth
controls the cycle rank of their supports, while the preceding finite
compression permits operator-valued entries.  The next two lemmas give the
support-rank bound and the walk count, including supports with cycles.

\begin{lemma}[Length versus cycle rank]
\label{lem:moment-support-rank}
Let S be a finite connected metric graph of total length at most Kg and girth
greater than g.  Then
\[
 \dim H_1(S;\mathbb Q)\leq \max\{1,\lceil20K\rceil\}^2.
\]
\end{lemma}

\begin{proof}
Set \(a=g/10\), and choose a maximal a-separated set P in the intrinsic
metric.  It is an a-net.  If it has at least two points, the open balls of
radius \(a/2\) around them are disjoint and each contains a path interval of
length \(a/2\) pointing toward another net point.  Hence
\(\abs P a/2\leq Kg\), and
\(\abs P\leq\max\{1,\lceil20K\rceil\}\), also covering a singleton net.
Join each pair of distinct net points at distance at most \(3a\) by one edge
of an abstract graph L, mapped to a chosen shortest path in S.  This map is
surjective on first homology.  To check this, subdivide a path of S into pieces
of length at most a, choose net points within a of the subdivision points,
and replace each piece and its two endpoint connectors by the chosen edge
path.  The replacement bounds a closed path of length at most \(6a<g\), which
is null-homotopic: cancellation of immediate reversals in a nontrivial closed
graph path leaves a path containing a simple cycle no longer than itself.
Connectors cancel consecutively around a loop.  Thus L is connected and its
mapped loops span first homology.  Finally
\(\dim H_1(S;\mathbb Q)\leq\dim H_1(L;\mathbb Q)\leq\abs P^2\).
\end{proof}

\begin{lemma}[Repeated-edge walk count]
\label{lem:repeated-edge-count}
Let a graph have maximum degree \(\Delta\geq1\).  Consider length-\(2h\)
closed walks from a fixed vertex, with no unoriented edge occurring exactly
once.  Multiple edges are treated as distinct edges.
If every support under consideration has cycle rank at most R, their
number is at most
\[
 [72(2R+2)^2\Delta]^h.
\]
\end{lemma}

\begin{proof}
A support has at most h edges.  Prune leaves other than the starting vertex,
retaining its rooted core K.  A tree leaves just that vertex.  Otherwise K
has no degree-one vertex except possibly the root and has the same cycle
rank as the support.  The identity
\[
 \sum_{v\in K}(\deg_K(v)-2)=2\dim H_1(K;\mathbb Q)-2
\]
gives \(\deg_K(v)\leq2R+2\).  The removed pieces are trees attached at single
core vertices, since neither a cycle nor a path between two core vertices
could be removed by this pruning.

There are at most \(4^e\Delta^e\) possible rooted connected subgraphs K with
e edges.  For clarity this includes cyclic K.  Fix orders on ambient incident
edges and explore K by \emph{marking edges}.  When an unmarked K-edge is
available, mark and traverse it, push it onto a stack, and recursively explore
at its endpoint, even if that vertex was already visited.  When the recursive
call has exhausted its unmarked edges, return along the pushed edge and pop
it.  Every edge is pushed and popped exactly once.  The balanced push/pop word
has at most \(4^e\) possibilities, and at each push an ambient edge has at
most \(\Delta\) choices.  These data recover the traversal and its edge set.

For fixed K, classify each walk step as a core step, a tree descent away from
the core, or a forced tree return.  There are at most \(3^{2h}\) type words.
Each core step has at most \(2R+2\) choices.  If t denotes the number of
descents and c the number of core steps, closure gives \(2t+c=2h\).  Every
core edge is traversed at least twice, so \(c\geq2e\) and \(t\leq h-e\).
Descents have at most \(\Delta\) choices and returns are determined by their
stack.  Thus the count is bounded by
\[
 \sum_{e=0}^h4^e\Delta^e3^{2h}(2R+2)^{2h}\Delta^{h-e}
 \leq[72(2R+2)^2\Delta]^h,
\]
using \(h+1\leq2^h\).  Invalid formal encodings only increase this upper
bound.  No assertion that a short backtracking walk is a simple path is used.
\end{proof}

\begin{lemma}[Centered block concentration]
\label{lem:centered-voltage-block}
Fix a surviving incidence type and one of its same or adjacent layer blocks.
Assign independent samples from Lemma~\ref{lem:product-sampler} to its
unoriented edges.  Let \(X_{ik}\) be its unnormalized centered convolution
matrix.  For every \(\eta>0\), the choices can be arranged in the order
described below so that
\begin{equation}
\label{eq:centered-block-tail}
 \Pr\{\norm{X_{ik}}>\eta q_{\tau,ik}\}\leq2^{-6H_i},
\end{equation}
uniformly in height.  All estimates are under the independent sampling law,
without conditioning on voltage girth.
\end{lemma}

\begin{proof}
An edge entry is \(W_e\lambda(v_e)-\mathbb E[W_e\lambda(v_e)]\) and has norm
at most \(2M_\eta\).  Double the block off the diagonal to make it
selfadjoint.  Its absolute row and column coefficient sums are at most
\(C_*M_\eta D\).

We first make precise why a fixed finite-test exponent is available before
choosing s.  Require \(H_0\geq s\), so
\[
 D\leq2^{.60005H_i},\qquad
 \abs{\mathscr A_i}=2^{H_i+.2s}\leq2^{1.2H_i}.
\]
After raising \(H_0\) to absorb fixed constants, the row bound is at most
\(2^{2H_i}\), and the number of base and matrix indices is at most
\(2^{3H_i}\).  The coordinates of every possible voltage and inverse have
the fixed exponent of Lemma~\ref{lem:product-sampler}.  Apply
Lemma~\ref{lem:heisenberg-folner-test} with a boundary exponent larger than
four.  Its box error tends to zero exponentially even after the row loss, and
its total finite compression dimension is at most \(2^{K_1H_i}\) for a fixed
\(K_1\).  Enlarge \(H_0\) further so the box error is at most
\(\eta q_{\tau,ik}/2\).  These enlargements do not change \(K_1\).

Now fix \(C>2(K_1+6)\), increasing it strictly if needed, and put
\[
 h=\left\lceil\frac{CH_i}{\log_2d}\right\rceil.
\]
Every no-singleton support in the \(2h\)-moment has at most h base edges.
It lies in the two fixed layers of this block.  Their base girth is at least
\(.9H_i/\log_2q\) after increasing \(H_0\), since adjacent H-values differ
by only b.  Because \(\log_2q/\log_2d<1.25\), the ratio of support length
to this girth is at most \(2C\), again increasing \(H_0\) for the ceiling.
Lemma~\ref{lem:moment-support-rank} gives a rank bound depending on C alone.
This is a use of \emph{base} girth, not the as-yet unproved voltage girth.

Let \(\widehat X\) denote the finite selfadjoint compression.  Expand
\(\mathbb E\operatorname{Tr}(\widehat X^{2h})\) along base vertices and
individual edges.  Any summand with a singleton unoriented edge has expectation
zero: after conditioning on the other edge draws it is a fixed left factor,
one centered compressed entry or its adjoint, and a fixed right factor.
Compression does not change this independence.  For every other summand its
absolute trace on the deck and matrix indices is at most their dimension
times \((2M_\eta)^{2h}\).  Lemma~\ref{lem:repeated-edge-count}, with
\(\Delta\leq C_*D\), proves
\[
 \mathbb E\operatorname{Tr}(\widehat X^{2h})
 \leq2^{K_1H_i}(C_1D)^h,
\]
where \(C_1\) is fixed after C and the coefficient data, but before s.
Since \(\widehat X\) is selfadjoint, its largest absolute eigenvalue to the
power \(2h\) is at most this nonnegative trace.  Markov's inequality and
\eqref{eq:analysis-degree-bounds} give
\[
 \Pr\{\norm{X_{ik}}>\eta q_{\tau,ik}\}
 \leq2^{K_1H_i}\left(\frac{C_2}{\eta^2d}\right)^h.
\]
Choose s large enough that \(C_2/\eta^2\leq d^{1/2}\).  The last expression
is at most \(2^{(K_1-C/2)H_i}\leq2^{-6H_i}\), by the choice of C.
The potentially large dependence of \(C_2\) on C is not circular: s is
chosen after both constants.  This proves the claim.
\end{proof}

\subsection{Simultaneous selection over all heights}
\label{subsec:analysis-simultaneous}

\begin{proof}[Proof of Proposition~\ref{prop:sampled-symbols}]
For a block put \(a_{\tau,ik}=q/q_{\tau,ik}\).  This scalar lies in a fixed
compact subinterval of \((0,\infty)\).  Apply
Lemma~\ref{lem:product-sampler}, with smaller tolerances, to the symbol
\(a_{\tau,ik}S_{\tau,ik}\).  Multiplication of the profile family by these
bounded scalars preserves every uniform estimate.  Round these already
rescaled weights to one finite net, and normalize the weighted incidence by
the common denominator q.  Within a block use the same law at every edge,
independently over all underlying unoriented edges.  Its expectation is
\[
 \frac{q_{\tau,ik}}q\,U_{\tau,ik}\otimes
       \mathbb E[W_e\rho(h_e)],
\]
whose distance from \(U_{\tau,ik}\otimes S_{\tau,ik}\) is bounded by a fixed
constant times the assigned one-norm error.
Lemma~\ref{lem:centered-voltage-block}, with its tolerance reduced using
\(q_{\tau,ik}/q\leq C_*\), supplies the remaining error.  This common-q
normalization avoids any need to recover an unbounded height from a label.

We verify the voltage-girth probability for the same law.  A forbidden closed
lift projects to a cyclically nonbacktracking closed base walk.  Test all such
base walks of length at most \(\lceil6.03m_i\rceil\) from height i.  For large
\(m_i\), this length is below \(6.04m_i\); hence the local bound in
Proposition~\ref{prop:pruning} says that each is a non-law.  This also removes
any rounding issue at exactly
\(6.03m_i\).  Along such a walk,
\[
 H_{\min}\geq H_i-b\lceil6.03H_i/\log_2q\rceil
 \geq(1-7b/\log_2q)H_i
\]
after raising the initial height.  Taking \(b/\log_2q\leq10^{-6}\),
Lemma~\ref{lem:polynomial-avoidance} bounds its identity probability by
\(2^{-40H_i}\), with all fixed factors absorbed by the initial height.

The number of possible rooted walk instructions is bounded by
\[
 (\abs{\mathscr A_i}+\abs{\mathscr B_i})
   \sum_{a\leq\lceil6.03m_i\rceil}(C_*D)^a
 \leq2^{9H_i}
\]
for sufficiently large s and then \(H_0/s\).  Indeed its leading exponent
divided by \(H_i\) is
\[
 1+6.03\frac{.60005}{.50005}<8.24;
\]
the finite type factor, the side-size offset \(.2s\), and rounding fit the
remaining margin.  The union bound for height i is therefore at most
\(2^{-31H_i}\).  It counts every edge, even those with realized weight zero.
Triviality of holonomy is independent of the chosen point in the initial
deck fiber, so no additional union over that infinite fiber is needed.

There are finitely many block types at each height.  Since
\(H_i=H_0+bi\), the sum over all heights of their norm-failure bounds and the
girth-failure bounds can be made less than one by raising \(H_0\).  No
conditioning on the girth events has entered the centered moment calculation.
For completeness, one can obtain a simultaneous choice without invoking an
infinite-product measure.  Enumerate all the required block inequalities and
walk tests.  Any finite list involves finitely many edge variables and, by
the same summed probability bound, has a satisfying assignment.  Extend those
assignments arbitrarily to all other edges.  Each edge has finitely many
possible voltage/net-weight pairs, so successive subsequences stabilize the
choice at the first edge, the second edge, and so on.  The diagonal limit
satisfies every condition, because each condition depends on finitely many
edge choices.  Fix such an assignment.

Finally append the finite matrix-weight index to each old label, and use the
inverse label with adjoint weight on its reverse.  Each refined positive
label now specifies its matrix weight, and every type uses the common
normalizer \(q^{-1}\).  These are the coefficients of the finite Cayley
polynomial in Equation~\eqref{eq:finite-cayley-polynomial}.  Its comparison
with the sampled blocks on retained sheet domains, and the estimate on the
remaining edges, are proved in Section~\ref{sec:assembly}.
Equality of refined words implies equality of
the original words, so all original uniqueness conclusions persist.  The
original leading label count is \(d2^r=2^{1.00005s}\); prescribed types, the
weight net and extras contribute height-independent factors, and the residue
flags contribute \(2^{O(\log s)}\).  Large voltage coordinates are separate
edge data and do not enter this alphabet.  Since
\(\log_2(q^2)=1.0001s\), the displayed sharp margin follows.  The ratio is
more precisely \(2^{-.000025s+O(\log s)+O_{b,\eta,\mathrm{profiles}}(1)}\),
not a claimed constant-prefactor exact decay.  This proves all assertions.
\end{proof}

\subsection{Choice order and the interface with constant extraction}

To apply Proposition~\ref{prop:sampled-symbols} along a sequence whose errors
tend to zero, fix b large enough for the uniform gap \(\rho<1\) of the
unweighted incidences, and keep b fixed throughout.  Also fix \(\alpha\),
the Bott profile, enlarged support cutoffs, and the matrix sizes.  These
choices give uniform profile bounds and homotopy moduli throughout the
accuracy sequence.

At accuracy \(\epsilon\), choose a finite extraction length k sufficiently
large for the gap term in Equation~\eqref{eq:constant-extraction-error}.
Choose a finite mesh of the
two adjacent-coordinate homotopies using their uniform profile modulus.
Collect the cutoffs, mesh symbols, parity selections and initial selection
into the finite list in Proposition~\ref{prop:sampled-symbols}.  Set their
individual tolerances small enough compared with \(\epsilon/k\), and choose
the finite Fourier boxes and weight net.  Then fix the finite-test exponent
\(K_1\), the moment constant C, choose s sufficiently large for all
concentration, label and upstream graph requirements, and finally choose
\(H_0\) sufficiently large.  In abbreviated form the order is
\[
 b,\ \text{profiles and matrix size};\qquad
 \epsilon,k,\text{mesh, types, net};\qquad
 K_1,C,s,H_0.
\]
The extraction length and finite type count may grow along the accuracy
sequence; they are fixed inside each tower before s is chosen.

For the assembly interface, let \(\chi_i\) be one of the larger cutoffs equal
to one on every required incident symbol support.  The preceding proposition
supplies, in particular, the full-norm approximations to
\[
 U_{ii}\otimes T_i\chi_iT_i^*,\qquad
 U_{ik}\otimes T_if_{ik}T_k^*,\quad\abs{i-k}\leq1.
\]
The constant-extraction argument in Section~\ref{sec:assembly} uses finite
products of these same-layer blocks and their adjoints, followed by a desired
cross block.  The cutoffs leave the target compact
symbols unchanged by Lemma~\ref{lem:supported-spreading}.  Proposition~
\ref{prop:sampled-symbols} controls each factor on the whole regular space,
including the complementary modes.  For a product of r bounded factors,
replacing each factor by an approximant of error \(\delta\) costs at most
\(r\delta\) times the appropriate fixed product norm bound, by telescoping.

Because b is fixed, finite meshes approximate the model homotopies uniformly
in both height and construction index.  Passing from these abstract block
estimates to operators in the physical reduced algebra requires the domain
and residual-forest conclusions of
Proposition~\ref{prop:domains} and the weighted forest estimate in the
proof of Lemma~\ref{lem:finite-profile-realization}.  Once the finite mesh
operators have been realized there, linear interpolation in each coordinate
reduced algebra has error tending to zero uniformly in the path parameter.
These uniform approximations give the continuous paths in the norm quotient
established in Lemma~\ref{lem:source-reduced-membership}.
